\chapter{Variables aléatoires continue}
















%=====================================================================
\section{Loi uniforme (continue)}
%=====================================================================

Analogue continu de la loi uniforme discrète~: c'est cette fois la
\emph{densité} qui est constante sur un intervalle.

\medskip

\begin{mdframed}
\begin{definition}[Loi uniforme continue]
Soit $I\subset\mathbb{R}$ un intervalle borné, de bornes $a=\inf I$ et $b=\sup I$
($a<b$), de longueur $|I|=b-a$. On dit que $X$ \emph{suit la loi uniforme sur
$I$}, et l'on note $X\sim\mathcal{U}(a,b)$, si $X$ admet la densité
\[
  f_X(x)=\frac{1}{b-a}\,\mathbb{1}_{[a,b]}(x).
\]
\end{definition}
\end{mdframed}

\medskip

\begin{mdframed}
\begin{proposition}[Fonction de répartition]
Pour $X\sim\mathcal{U}(a,b)$,
\[
  F_X(t)=
  \begin{cases}
    0, & t<a,\\[2pt]
    \dfrac{t-a}{b-a}, & a\le t\le b,\\[6pt]
    1, & t>b.
  \end{cases}
\]
La notation $\mathcal{U}(a,b)$ ne précise pas si les bornes sont incluses~: cela n'a
aucune incidence, un point étant de probabilité nulle.
\end{proposition}
\end{mdframed}

\medskip

\begin{mdframed}
\begin{proposition}[Espérance, variance, moments]
\label{prop:unifc-moments}
Si $X\sim\mathcal{U}(a,b)$, alors

\begin{itemize}
\item $\mathbb{E}[X]=\frac{a+b}{2}$,
\item $\mathbb{V}(X)=\frac{(b-a)^2}{12}$,
\item $M_X(t)=\frac{e^{tb}-e^{ta}}{t(b-a)}$ avec $t\ne0$.
\item Plus généralement $\displaystyle\mathbb{E}[X^n]=\frac{b^{n+1}-a^{n+1}}{(n+1)(b-a)}$.
\end{itemize}
\end{proposition}
\end{mdframed}

\begin{proof}
$$
\mathbb{E}[X]=\frac1{b-a}\int_a^b x\,dx=\frac{b^2-a^2}{2(b-a)}=\frac{a+b}{2}
$$ 

et
$$
\mathbb{E}[X^2]=\frac1{b-a}\int_a^b x^2\,dx=\frac{b^3-a^3}{3(b-a)}
=\frac{a^2+ab+b^2}{3}. 
$$

D'où

\begin{align*}
  \mathbb{V}(X)&=\frac{a^2+ab+b^2}{3}-\frac{(a+b)^2}{4} \\
  &=\frac{4(a^2+ab+b^2)-3(a+b)^2}{12} \\
&=\frac{(b-a)^2}{12}.
\end{align*}

Enfin $M_X(t)=\frac1{b-a}\int_a^b e^{tx}\,dx=\frac{e^{tb}-e^{ta}}{t(b-a)}$.
\end{proof}

\begin{remarque}[L'intervalle est nécessairement borné]
On ne peut pas prendre $a=-\infty$ ou $b=+\infty$~: il n'existe pas de loi
uniforme sur $\mathbb{R}$ tout entier (de même qu'il n'existe pas de loi uniforme sur
$\mathbb{N}$).
\end{remarque}

\begin{exercice}
\label{ex:unif-affine}
Soit $X\sim\mathcal{U}(0,1)$ et $a,b\in\mathbb{R}$ avec $a\ne0$. Déterminer la loi de
$aX+b$. \emph{(Réponse~: $\mathcal{U}$ des bornes $b$ et $a+b$.)}
\end{exercice}

\newpage
%=====================================================================
\section{Loi normale (ou gaussienne)}
%=====================================================================

\medskip

\begin{mdframed}
\begin{definition}[Loi normale]
Soient $\mu\in\mathbb{R}$ et $\sigma^2>0$. On dit que $X$ \emph{suit la loi normale
(ou gaussienne) de moyenne $\mu$ et variance $\sigma^2$}, et l'on note
$X\sim\mathcal N(\mu,\sigma^2)$, si $X$ admet la densité
\[
  f_X(x)=\frac{1}{\sqrt{2\pi\sigma^2}}\;\exp\!\left(-\frac{(x-\mu)^2}{2\sigma^2}\right).
\]
Si $\mu=0$, $X$ est dite \emph{centrée}~; si de plus $\sigma^2=1$, \emph{réduite}.
\end{definition}
\end{mdframed}

\medskip

\begin{mdframed}
\begin{proposition}[$f_X$ est bien une densité]
On a $\int_\mathbb{R} f_X=1$. Par le changement de variable $t=(x-\mu)/(\sigma\sqrt2)$
on se ramène à l'intégrale de Gauss $\int_\mathbb{R} e^{-t^2}\,dt=\sqrt\pi$.
\end{proposition}
\end{mdframed}

\medskip

\begin{mdframed}
\begin{notation}[Fonction $\Phi$]
La fonction de répartition de la loi \emph{centrée réduite} $\mathcal N(0,1)$
est notée
\[
  \Phi(t)=\int_{-\infty}^{t}\frac{1}{\sqrt{2\pi}}\,e^{-y^2/2}\,dy.
\]
Elle n'a pas d'expression élémentaire (elle est reliée à la fonction d'erreur
$\operatorname{erf}(x)=\frac{2}{\sqrt\pi}\int_0^x e^{-y^2}\,dy$), mais elle est tabulée et
disponible dans tout logiciel de calcul.
\end{notation}
\end{mdframed}

\medskip

\begin{mdframed}
\begin{lemme}[Transformation affine --- \textnormal{à retenir}]
\label{lem:normale-affine}
Si $N\sim\mathcal N(0,1)$ et $\mu,\sigma\in\mathbb{R}$ avec $\sigma\ne0$, alors
$\mu+\sigma N\sim\mathcal N(\mu,\sigma^2)$.
\end{lemme}
\end{mdframed}

\begin{proof}
Cas $\sigma>0$. Pour $t\in\mathbb{R}$,

\begin{align*}
  \mathbb{P}(\mu+\sigma N\le t) &=\mathbb{P}\!\Big(N\le\tfrac{t-\mu}{\sigma}\Big) \\
  &=\int_{-\infty}^{(t-\mu)/\sigma}\frac{e^{-y^2/2}}{\sqrt{2\pi}}\,dy \\
  &=\int_{-\infty}^{t}\frac{1}{\sqrt{2\pi\sigma^2}}\,
  e^{-(z-\mu)^2/(2\sigma^2)}\,dz,
\end{align*}

après le changement de variable affine $y=(z-\mu)/\sigma$. On reconnaît
$F_{\mathcal N(\mu,\sigma^2)}$. Pour $\sigma<0$, on note que $N$ et $-N$ ont
même loi, donc $\mu+\sigma N$ et $\mu-\sigma N$ aussi.
\end{proof}

\medskip

\begin{mdframed}
\begin{proposition}[Espérance, variance, moments]
Si $X\sim\mathcal N(\mu,\sigma^2)$, alors

\begin{itemize}
\item $\mathbb{E}[X]=\mu$, 
\item $\mathbb{V}(X)=\sigma^2$,
\item $M_X(t)=\exp\!\Big(\mu t+\tfrac12\sigma^2 t^2\Big)$, avec $t\in\mathbb{R}$.
\end{itemize}

Les moments centrés sont, pour $N\sim\mathcal N(0,1)$:

\begin{itemize}
\item $\mathbb{E}[N^{2k+1}]=0$ 
\item $\mathbb{E}[N^{2k}]=(2k-1)!!=\dfrac{(2k)!}{2^k\,k!}$.
\end{itemize}
\end{proposition}
\end{mdframed}

\begin{proof}
Écrivons $X=\mu+\sigma N$. Par imparité, $\mathbb{E}[N]=0$; 

une intégration par parties donne $\mathbb{E}[N^2]=\int y^2\frac{e^{-y^2/2}}{\sqrt{2\pi}}\,dy=1$. 

Donc $\mathbb{E}[X]=\mu$ et $\mathbb{V}(X)=\sigma^2\mathbb{V}(N)=\sigma^2$. Pour la FGM,

\[
  M_N(t)=\int_\mathbb{R} e^{ty}\frac{e^{-y^2/2}}{\sqrt{2\pi}}\,dy
  =e^{t^2/2}\int_\mathbb{R}\frac{e^{-(y-t)^2/2}}{\sqrt{2\pi}}\,dy=e^{t^2/2},
\]

d'où $M_X(t)=e^{\mu t}M_N(\sigma t)=e^{\mu t+\sigma^2t^2/2}$. Les moments de
$N$ s'obtiennent en identifiant le développement en série de $e^{t^2/2}$.
\end{proof}

\medskip

\begin{mdframed}
\begin{proposition}[Convention $\sigma=0$]
On pose $\mathcal N(\mu,0):=\delta_\mu$ (loi de Dirac en $\mu$) : ce n'est plus
une loi à densité.
\end{proposition}
\end{mdframed}

\begin{remarque}[Forme de la densité]
La densité est la « courbe en cloche » de Gauss, symétrique autour de $\mu$,
croissante puis décroissante~; plus $\sigma^2$ est petit, plus elle est
resserrée autour de $\mu$.
\end{remarque}

\begin{exercice}[Encadrement de la queue gaussienne]
\label{ex:queue-gauss}
Soit $N\sim\mathcal N(0,1)$. Montrer, par intégrations par parties successives
de $1-\mathbb{P}hi(t)$, l'encadrement (excellent pour $t$ grand)~:
\[
  \Big(\frac1t-\frac1{t^3}\Big)e^{-t^2/2}\le \sqrt{2\pi}\,\mathbb{P}(N\ge t)
  \le \frac1t\,e^{-t^2/2},\qquad t>0.
\]
\end{exercice}

\newpage












































































\subsection{Approximation gaussienne de la loi binomiale}

\medskip

\begin{mdframed}
\begin{theoreme}[de Moivre-Laplace]
\label{th:moivre-laplace}
Soit $p\in\,]0,1[$ et, pour tout $n$, $X_n\sim\mathcal{B}(n,p)$. Alors pour
tous $-\infty\le a<b\le+\infty$,
\[
  \mathbb{P}\!\Big(np+a\sqrt{np(1-p)}\le X_n\le np+b\sqrt{np(1-p)}\Big)
  \xrightarrow[n\to\infty]{}\ \Phi(b)-\Phi(a).
\]
\end{theoreme}
\end{mdframed}

\begin{remarque}[Heuristique]
Le théorème dit qu'à $p$ fixé et $n$ grand,
\[
  X_n\approx np+\sqrt{np(1-p)}\;N,\qquad N\sim\mathcal N(0,1),
\]
c'est-à-dire qu'on approche $\mathcal{B}(n,p)$ par
$\mathcal N\big(np,\,np(1-p)\big)$. Historiquement, de Moivre traita le cas
$p=1/2$ (1733), Laplace le cas général (1812)~; c'est un cas particulier du
théorème central limite. En pratique, l'approximation est jugée bonne dès que
$np(1-p)\ge10$.
\end{remarque}

\begin{remarque}[Correction de continuité]
Pour approcher $\mathbb{P}(a\le X_n\le b)$ avec $a,b$ \emph{entiers}, on approche
chaque atome par un intervalle de largeur $1$ centré en lui~:
$\mathbb{P}(X_n=k)\approx\mathbb{P}(k-\tfrac12<X\le k+\tfrac12)$ avec
$X\sim\mathcal N(np,np(1-p))$. En sommant sur $k\in\{a,\dots,b\}$ et en posant
$N\sim\mathcal N(0,1)$,
\[
  \mathbb{P}(a\le X_n\le b)\approx
  \Phi\!\left(\frac{b+\tfrac12-np}{\sqrt{np(1-p)}}\right)
  -\Phi\!\left(\frac{a-\tfrac12-np}{\sqrt{np(1-p)}}\right).
\]
Les demi-entiers $a-\tfrac12$ et $b+\tfrac12$ (au lieu de $a,b$) constituent
la \emph{correction de continuité}~; elle améliore sensiblement
l'approximation.
\end{remarque}

%=====================================================================
\section{Loi exponentielle}
%=====================================================================

Analogue continu de la loi géométrique.

\medskip

\begin{mdframed}
\begin{definition}[Loi exponentielle]
Soit $\lambda>0$. On dit que $X\sim\mathcal{E}(\lambda)$ si $X$ admet la densité
\[
  f_X(x)=\lambda\,e^{-\lambda x}\,\mathbb{1}_{\mathbb{R}_+}(x).
\]
\end{definition}
\end{mdframed}

\medskip

\begin{mdframed}
\begin{proposition}[Fonction de répartition et queue]
Pour $X\sim\mathcal{E}(\lambda)$,

$$
  F_X(t)=\big(1-e^{-\lambda t}\big)\,\mathbb{1}_{\mathbb{R}_+}(t),
$$

Et en particulier

$$
  \mathbb{P}(X>t)=e^{-\lambda t} \qquad (t\ge0).
$$
\end{proposition}
\end{mdframed}

\medskip

\begin{mdframed}
\begin{proposition}[Espérance, variance, moments]
Si $X\sim\mathcal{E}(\lambda)$, alors

\begin{itemize}
\item $\mathbb{E}[X]=\frac1\lambda$,
\item $\mathbb{V}(X)=\frac1{\lambda^2}$,
\item $\mathbb{E}[X^n]=\frac{n!}{\lambda^n}$,
\item $M_X(t)=\frac{\lambda}{\lambda-t}$ avec $t<\lambda$.
\end{itemize}
\end{proposition}
\end{mdframed}

\begin{proof}
$$
\mathbb{E}[X^n]=\int_0^\infty x^n\lambda e^{-\lambda x}\,dx=\frac{n!}{\lambda^n}
$$

(intégrale gamma). En particulier 

$$
\mathbb{E}[X]=1/\lambda, \qquad \mathbb{E}[X^2]=2/\lambda^2,
$$

d'où 

$$
\mathbb{V}(X)=2/\lambda^2-1/\lambda^2=1/\lambda^2. 
$$

Enfin

$$
M_X(t)=\int_0^\infty e^{tx}\lambda e^{-\lambda x}\,dx=\frac{\lambda}{\lambda-t}
$$

si $t<\lambda$.
\end{proof}

\medskip

\begin{mdframed}
\begin{proposition}[Absence de mémoire --- \textnormal{caractérisation à retenir}]
\label{prop:exp-memoire}
Soit $X$ une v.a.r. Alors $X\sim\mathcal{E}(\lambda)$ pour un certain $\lambda>0$ si
et seulement si
\[
  \mathbb{P}(X\in\mathbb{R}_+)=1
  \quad\text{et}\quad
  \mathbb{P}(X>t+s\mid X>t)=\mathbb{P}(X>s),\qquad \forall\,t,s\in\mathbb{R}_+.
\]
C'est l'exact analogue continu de la loi géométrique.
\end{proposition}
\end{mdframed}

\begin{proof}
La condition s'écrit $g(t+s)=g(t)g(s)$ avec $g(t)=\mathbb{P}(X>t)$. Pour $p,q\in\mathbb{N}^*$,

$$
g(p/q)=g(1/q)^p=g(1)^{p/q};
$$ 

comme $g$ est décroissante, on prolonge à tout $t\in\mathbb{R}_+$ : 

$$
g(t)=g(1)^t=e^{-\lambda t} \quad \text{avec} \quad \lambda=-\log g(1)\ge0.
$$ 

On reconnaît la queue de $\mathcal{E}(\lambda)$. La réciproque est un calcul direct:

$$
\frac{\mathbb{P}(X>t+s)}{\mathbb{P}(X>t)}=\frac{e^{-\lambda(t+s)}}{e^{-\lambda t}}=e^{-\lambda s}.
$$
\end{proof}

\medskip

\begin{mdframed}
\begin{lemme}[Changement d'échelle]
Si $X\sim\mathcal{E}(\lambda)$ et $\mu>0$, alors $\mu X\sim\mathcal{E}(\lambda/\mu)$.
\end{lemme}
\end{mdframed}

\begin{proof}
Pour $x\ge0$, 

$$
F_{\mu X}(x)=\mathbb{P}(X\le x/\mu)=1-e^{-(\lambda/\mu)x}, 
$$

qui est la fonction de répartition de $\mathcal{E}(\lambda/\mu)$.
\end{proof}

%=====================================================================
\section{Loi gamma}
%=====================================================================

\medskip

\begin{mdframed}
\begin{definition}[Loi gamma]
Soient $\alpha,\beta>0$. On dit que $X\sim\Gamma(\alpha,\beta)$ si $X$ admet
la densité
\[
  f_X(x)=\frac{\beta^\alpha}{\Gamma(\alpha)}\,x^{\alpha-1}e^{-\beta x}
  \,\mathbb{1}_{]0,+\infty[}(x),
\]
où $\displaystyle\Gamma(\alpha)=\int_0^\infty x^{\alpha-1}e^{-x}\,dx$ est la
fonction gamma d'Euler, qui vérifie $\Gamma(n)=(n-1)!$ pour $n\in\mathbb{N}^*$ et
interpole donc la factorielle.
\end{definition}
\end{mdframed}

\medskip

\begin{mdframed}
\begin{proposition}[Espérance, variance, moments]
\label{prop:gamma-moments}
Si $X\sim\Gamma(\alpha,\beta)$, alors

\begin{itemize}
\item $\mathbb{E}[X]=\frac{\alpha}{\beta}$,
\item $\mathbb{V}(X)=\frac{\alpha}{\beta^2}$,
\item $\mathbb{E}[X^n]=\frac{\Gamma(\alpha+n)}{\Gamma(\alpha)\,\beta^n} =\frac{\alpha(\alpha+1)\cdots(\alpha+n-1)}{\beta^n}$,
\item $M_X(t)=\big(1-t/\beta\big)^{-\alpha}$ pour $t<\beta$.
\end{itemize}

La fonction de répartition n'a pas de forme close simple.
\end{proposition}
\end{mdframed}

\begin{proof}
Par définition de $\Gamma$,
\[
  \mathbb{E}[X^n]=\frac{\beta^\alpha}{\Gamma(\alpha)}\int_0^\infty x^{\alpha+n-1}
  e^{-\beta x}\,dx
  =\frac{\beta^\alpha}{\Gamma(\alpha)}\cdot\frac{\Gamma(\alpha+n)}{\beta^{\alpha+n}}
  =\frac{\Gamma(\alpha+n)}{\Gamma(\alpha)\beta^n}.
\]
Pour $n=1,2$ on obtient $\mathbb{E}[X]=\alpha/\beta$ et
$\mathbb{E}[X^2]=\alpha(\alpha+1)/\beta^2$, d'où $\mathbb{V}(X)=\alpha/\beta^2$. Le calcul de
$M_X$ est identique en remplaçant $\beta$ par $\beta-t$.
\end{proof}

\begin{mdframed}
\begin{proposition}[Lien avec l'exponentielle]
On a $\Gamma(1,\beta)=\mathcal{E}(\beta)$. Plus généralement, si $\alpha=r\in\mathbb{N}^*$ et
$X_1,\dots,X_r$ sont indépendantes de loi $\mathcal{E}(\beta)$, alors
$X_1+\dots+X_r\sim\Gamma(r,\beta)$.
\end{proposition}
\end{mdframed}

\begin{remarque}
Les valeurs \emph{non entières} de $\alpha$ donnent un sens à des « sommes non
entières » de lois $\mathcal{E}(\beta)$, tout comme $\Gamma$ interpole la factorielle.
La loi gamma est aussi reliée à la loi de Poisson, d'où de nombreuses
applications (files d'attente, fiabilité).
\end{remarque}

%=====================================================================
\section{Loi de Cauchy}
%=====================================================================
\medskip

\begin{mdframed}
\begin{definition}[Loi de Cauchy]
Soit $a>0$. On dit que $X\sim\mathcal{C}(a)$ si $X$ admet la densité
\[
  f_X(x)=\frac{a}{\pi}\,\frac{1}{a^2+x^2},\qquad x\in\mathbb{R}.
\]
\end{definition}
\end{mdframed}

\medskip

\begin{mdframed}
\begin{proposition}[Fonction de répartition]
Pour $X\sim\mathcal{C}(a)$,
\[
  F_X(t)=\int_{-\infty}^t\frac a\pi\frac{dx}{a^2+x^2}
  =\frac12+\frac1\pi\arctan\!\Big(\frac ta\Big).
\]
\end{proposition}
\end{mdframed}

\medskip

\begin{mdframed}
\begin{proposition}[Aucun moment n'existe]
\label{prop:cauchy-moments}
Si $X\sim\mathcal{C}(a)$, alors $\mathbb{E}[|X|]=+\infty$~: l'espérance \emph{n'existe pas}
(et \emph{a fortiori} ni la variance, ni la FGM). La fonction caractéristique
vaut néanmoins $\mathbb{E}[e^{itX}]=e^{-a|t|}$.
\end{proposition}
\end{mdframed}

\begin{proof}
$$
\displaystyle\int_\mathbb{R} |x|\,\frac a\pi\frac{dx}{a^2+x^2}
=\frac{2a}\pi\int_0^\infty\frac{x}{a^2+x^2}\,dx
$$

diverge (l'intégrande se comporte en $1/x$ à l'infini). Donc $X$ n'est pas intégrable.
\end{proof}

\begin{remarque}
La densité de Cauchy a elle aussi une allure « en cloche », mais à
\emph{queues lourdes}~: elle décroît en $1/x^2$, bien plus lentement que la
gaussienne, ce qui explique la non-existence des moments.
\end{remarque}

\begin{exercice}[Loi de l'inverse]

\begin{enumerate}
  \item Montrer que $\arctan(t)+\arctan(1/t)=\dfrac\pi2\big(\mathbb{1}_{t>0}-\mathbb{1}_{t<0}\big)$
        pour tout $t\ne0$.
  \item Soit $X\sim\mathcal{C}(a)$ (on a $\mathbb{P}(X\ne0)=1$). En déduire la loi de
        $1/X$. \emph{(Réponse~: $1/X\sim\mathcal{C}(1/a)$.)}
\end{enumerate}
\end{exercice}

\newpage
%=====================================================================
\section{Autres lois à densité}
%=====================================================================

Comme dans le cas discret, ces lois ne sont pas à connaître par c\oe ur.

\subsection{Loi de Weibull}

\medskip

\begin{mdframed}
\begin{definition}[Loi de Weibull]
Soient $\alpha,\beta>0$. On dit que $X\sim\mathrm{Weibull}(\alpha,\beta)$ si sa fonction de
répartition est
\[
  F_X(x)=
  \begin{cases}
    1-\exp\!\big(-(x/\alpha)^\beta\big), & x\ge0,\\[2pt]
    0, & x<0,
  \end{cases}
\]
soit, en dérivant, la densité
$f_X(x)=\dfrac\beta\alpha\Big(\dfrac x\alpha\Big)^{\beta-1}
\exp\!\big(-(x/\alpha)^\beta\big)\mathbb{1}_{x\ge0}$.
\end{definition}
\end{mdframed}

\medskip

\begin{mdframed}
\begin{proposition}[Premiers moments]
Si $X\sim\mathrm{Weibull}(\alpha,\beta)$, alors
\[
  \mathbb{E}[X]=\alpha\,\Gamma\!\Big(1+\tfrac1\beta\Big),\qquad
  \mathbb{V}(X)=\alpha^2\!\left[\Gamma\!\Big(1+\tfrac2\beta\Big)
  -\Gamma\!\Big(1+\tfrac1\beta\Big)^2\right].
\]
\end{proposition}
\end{mdframed}

\begin{remarque}
La loi de Weibull généralise l'exponentielle~: pour $\beta=1$ on retrouve
$\mathcal{E}(1/\alpha)$. Selon que $\beta<1$, $\beta=1$ ou $\beta>1$, le « taux de
panne » décroît, reste constant ou croît avec le temps~: d'où son usage en
modélisation de l'usure de pièces et en fiabilité.
\end{remarque}

\end{document}

\begin{mdframed}
A
\end{mdframed}